\documentclass[preprint,12pt]{elsarticle}

\usepackage{amsmath,amssymb}
\usepackage{graphicx}
\usepackage{hyperref}

\newcommand{\SU}{\mathrm{SU}}
\newcommand{\UU}{\mathrm{U}}
\newcommand{\Nf}{N_{\!f}}
\newcommand{\Nc}{N_{\!c}}
\newcommand{\Tr}{\mathrm{Tr}}
\newcommand{\adj}{\mathrm{adj}}
\newcommand{\MeV}{\,\mathrm{MeV}}
\newcommand{\GeV}{\,\mathrm{GeV}}
\newcommand{\TeV}{\,\mathrm{TeV}}

\begin{document}

\begin{frontmatter}

\title{Quadratic charge--mass relation from gauged family symmetry}

\author{XXXXX}
\address{XXXXX}

\begin{abstract}
In any gauge theory, the classical field energy of a static source
scales as the square of its charge:
Gauss's law gives $F \propto q$, and the gauge kinetic term
$\mathcal{L} \supset F^2$ gives energy~$\propto q^2$.
Applied to a $\UU(3)$ family symmetry with charges
$q_i = z_0 + z_i$, this motivates the mass ansatz
$m_i = \mu\,q_i^2$, whose ratio
$K \equiv \sum q_i^2/(\sum q_i)^2 = 2/3$
is an algebraic identity of the charge normalisation.
We show that $n = 2$ is the unique nontrivial power
for which this ratio is independent of the charge direction
in the Cartan subalgebra
(proven analytically for integer~$n$;
verified numerically for real~$n$): the binomial coefficient
$n(n{-}1)(n{-}2)/6$ is nonzero for all real $n \neq 0,1,2$,
introducing $\alpha$-dependent terms through
$\sum z_i^3 \propto \cos 3\alpha$.
In a minimal $\mathcal{N}=1$ supersymmetric $\SU(3)$ gauge theory
with $\Nf = \Nc = 3$, the confined meson sector and its adjugate
preserve two holomorphic mass spectra satisfying
$K = 2/3$ and $K^{-1} = 2/3$ respectively,
where the $q^2$ dependence is an input hypothesis
protected by the non-renormalisation theorem.
Applied to the charged-lepton spectrum
(one-parameter prediction:
$m_\tau^{\rm pred} = 1776.97\MeV$,
$0.9\sigma$ from experiment)
and the down-quark spectrum
(zero-parameter phenomenological test: $1.2\sigma$),
the ratios agree with $K = 2/3$
to $0.001\%$ and $0.3\%$ respectively.
\end{abstract}

\end{frontmatter}

%----------------------------------------------------------------------
\section{Introduction}
The fermion masses and mixing parameters account for
13 of the 19 free parameters of the Standard Model
(9 masses, 3 CKM angles, 1 CP phase),
yet no principle within the SM constrains their values.
The empirical observation by Koide~\cite{Koide:1983qe}
that the charged-lepton masses satisfy
$(\sum m_i)/(\sum\sqrt{m_i})^2 = 2/3$ to $0.001\%$
has resisted a compelling theoretical explanation
for four decades, despite extensive work on
flavour symmetries~\cite{FroggattNielsen:1979},
composite models~\cite{Rivero:2005},
and discrete family groups~\cite{Ma:2006}.
We show that gauging a $\UU(3)$ family symmetry with the
standard gauge kinetic term motivates a quadratic
mass--charge relation,
$m_i = \mu\,q_i^2$,
and that if this ansatz is adopted as input,
$\mathcal{N}{=}1$ supersymmetric confinement
\emph{preserves} it to the infrared and simultaneously
generates an inverse sector, linking two
distinct fermion spectra through
the meson--adjugate duality of the confined phase.

The argument proceeds in three steps.
First, we present a \emph{field-energy argument}:
in any $\UU(1)$ gauge theory, the classical field energy
of a static charged source scales as~$q^2$.
Second, we show a \emph{uniqueness result}:
$n = 2$ is the only monomial power for which the
ratio $K = \sum q_i^n / (\sum q_i^{n/2})^2$ is independent
of the charge direction in the family Cartan subalgebra.
Third, we exhibit a minimal $\SU(3)$ SQCD model
whose confined phase preserves $K = 2/3$ (direct)
and simultaneously generates $K^{-1} = 2/3$ (inverse)
from a single gauge sector.

%----------------------------------------------------------------------
\section{Field-energy argument}

In any theory containing a $\UU(1)$ gauge field
with kinetic term
$\mathcal{L} \supset -(1/4g^2)F_{\mu\nu}^2$,
the classical field energy of a static source of charge~$q$
scales as~$q^2$.
Gauss's law is linear: $\nabla\!\cdot\!\mathbf{E} = g^2\rho$,
with $\rho_i(\mathbf{x}) = q_i\,\rho_0(\mathbf{x})$
for a source whose spatial profile~$\rho_0$ is
charge-independent.  Therefore
$\mathbf{E}_i(\mathbf{x}) = q_i\,\mathbf{E}_0(\mathbf{x})$
by linearity, and the stored field energy is
\begin{equation}
  U_i = \frac{1}{2g^2}\!\int\!d^3x\;\mathbf{E}_i^2
  \;=\; q_i^2\;\mu\,,
  \label{eq:field_energy}
\end{equation}
with $\mu = (2g^2)^{-1}\!\int\!d^3x\,\mathbf{E}_0^2$
generation-independent.

The result is universal: the quadratic gauge kinetic term
is the leading low-energy action for any abelian gauge field.
It is the gauge-theoretic version of the capacitor formula
$E = Q^2/(2C)$.
The integral in Eq.~(\ref{eq:field_energy}) is formally
UV-divergent for a point source;
confinement of the family gauge group provides the
finite size and renders~$\mu$ a finite function of
$\Lambda_F$ and~$g_F$, determined in principle by
the confining dynamics~\cite{Sumino:2009}.

We emphasise that the extension from classical field
energy to the mass of a confined composite is a
\emph{hypothesis}: the $q^2$ scaling is rigorous for
a static abelian source in the weak-field regime.
In a non-abelian confining theory, string tensions
depend on the full representation, not just
the Cartan charges; however, since the family
charges $q_i$ are Cartan eigenvalues of a single
representation, the abelian scaling
$q_i^2$ is the natural leading-order result.
We conjecture that it survives in the strongly-coupled
confining dynamics of the family gauge group.

The same $q^2$ structure persists in $\mathcal{N}{=}1$
supersymmetry: the bosonic Hamiltonian after $D$-term
elimination contains $\mathbf{E}^2/(2g_F^2)$ and
$g_F^2(\sum_i q_i|\phi_i|^2)^2/2$, again positive
forms sourced by~$q_i$.
This does not add probative force beyond the
classical argument---it merely confirms that SUSY
does not alter the $q^2$ scaling of the gauge
kinetic contribution.
In the confining phase, the total composite mass is
conjectured to receive a generation-dependent
piece $\Delta M_i = q_i^2\,\mu$, with $\mu$ set by
the non-perturbative dynamics.
If this is the dominant generation-dependent contribution
to the composite mass---i.e.\ any charge-independent
binding energy is universal---then
$m_i \approx m_0 + q_i^2\,\mu$.
Since $K$ involves $\sqrt{m_i}$, a nonzero $m_0$
would spoil $K = 2/3$; we therefore require
$m_0 \ll q_{\rm min}^2\,\mu \sim m_e$.
In the supersymmetric model of Section~5,
this is natural: the charge-independent part
of the superpotential is a constant
(the cosmological-constant term),
which vanishes in the SUSY limit and is
protected by the non-renormalisation theorem.
Physical SUSY breaking introduces
$m_0 \sim F/\Lambda$, negligible for
$F \ll m_e\,\Lambda$.
This is not the BPS mass $M \geq |Z|$
(linear in charge), which applies to
point-like dyons on the Coulomb branch.
Our composites are extended objects on the
confining branch: their mass receives
contributions from the field energy stored
in the confining flux tube, which has a
charge-independent spatial profile set
by~$\Lambda$ and~$g_F$.
The $q^2$ scaling is therefore a consequence
of the gauge kinetic term, not of the
central-charge algebra.

%----------------------------------------------------------------------
\section{Family charges}
Consider a $\UU(3)$ family symmetry.
The Cartan subalgebra is parametrised by a traceless
matrix $T = \mathrm{diag}(z_1, z_2, z_3)$ with
$\sum z_i = 0$ and a universal component $z_0 = 1/\sqrt{6}$
fixed by canonical normalisation
(see e.g.~\cite{Georgi:1999,Foot:1994}).
The total family charge of the $i$-th species is
\begin{equation}
  q_i \;=\; z_0 + z_i
  \;=\; \frac{1}{\sqrt{6}}
    + \frac{1}{\sqrt{3}}\cos\!\bigl(\alpha + \tfrac{2\pi(i{-}1)}{3}\bigr),
  \quad i = 1,2,3\,,
  \label{eq:qi}
\end{equation}
where $\alpha$ parametrises the direction in the
two-dimensional Cartan plane.
The normalisation conditions are
$\sum q_i = 3z_0 = \sqrt{3/2}$ and
$\sum q_i^2 = 3z_0^2 + \sum z_i^2 = 1/2 + 1/2 = 1$
(using $\sum z_i^2 = (1/3)\sum\cos^2(\alpha{+}2\pi i/3) = 1/2$).
Positivity of all three charges requires
$|z_i| < z_0$ for each~$i$, i.e.\
$\alpha$ is bounded to the interval in which
no $q_i$ vanishes; the physical value
$\alpha_{\rm lep} = 0.222$~rad lies well within
this domain.

%----------------------------------------------------------------------
\section{Why \texorpdfstring{$n=2$}{n=2}}
Since
$z_i = (1/\sqrt{3})\cos(\alpha + 2\pi i/3)$,
the power sums $\sum z_i^p$ vanish for $p = 1$
and are $\alpha$-independent for $p = 2$
(since $\sum \cos^2(\alpha + 2\pi i/3) = 3/2$
by the triple-angle identity), but
$\sum z_i^3 \propto \cos 3\alpha$ varies with~$\alpha$.
Now expand $\sum q_i^n = \sum(z_0+z_i)^n$
by the binomial theorem: the term
$\binom{n}{3}z_0^{n-3}\sum z_i^3$ is present
and $\alpha$-dependent for every $n \geq 3$.
For $n = 2$, only $p = 0,1,2$ contribute
and the numerator $\sum q_i^2 = 1$ is constant;
the denominator $(\sum q_i)^2 = 3/2$ is likewise
constant, giving
$K_2 = 2/3$ independent of~$\alpha$.
For any real $n \neq 0, 1, 2$, the generalised binomial
coefficient $n(n{-}1)(n{-}2)/6$ is nonzero,
so the numerator $\sum q_i^n$ acquires $\alpha$-dependent
terms through $\sum z_i^3 \propto \cos 3\alpha$.
For integer $n \geq 3$, this constitutes an
analytical proof that $K_n$ depends on~$\alpha$.
For non-integer $n$, the binomial expansion is an
infinite series; the leading Fourier mode
$\cos 3\alpha$ enters with nonzero coefficient, and
cancellation among the infinitely many higher-order
modes would require a fine-tuned identity.
We have verified numerically that no such
cancellation occurs for any tested value:
$K_n$ varies with~$\alpha$ for all sampled
real $n$ in $[-5, 10]$ at spacing $0.001$,
except~$n = 2$ (and the trivial $n = 0$);
we conjecture that no further exceptions exist.
For $n = 1$, the numerator is $\alpha$-independent
but the denominator $(\sum q_i^{1/2})^2$ is not.
For $n = 0$, $K_0 = 1/3$ trivially.
The argument generalises to $\UU(N)$: for
any~$N$, $\sum q_i^2$ is direction-independent
in the Cartan subalgebra
(since $\sum z_i^2 = \Tr(T^2)$ is a Casimir),
so $K_2 = 2/3$ holds for all~$N$.
The choice $N = 3$ is fixed by the
three observed generations.

%----------------------------------------------------------------------
\section{The minimal model: s-confinement}

We consider $\SU(3)$ supersymmetric QCD with
$\Nf = \Nc = 3$ flavours of chiral superfields:
fundamentals~$Q^i$ and antifundamentals~$\tilde{Q}_i$
($i = 1,2,3$), all neutral under the Standard Model.
We briefly review the infrared dynamics,
following~\cite{Seiberg:1994bz,IntriligatorSeibergShenker:1995,Csaki:1996zb};
the reader familiar with s-confinement may skip
to~Eq.~(\ref{eq:W}).

For $\Nf = \Nc$, the theory \emph{s-confines}:
unlike ordinary QCD, there is no chiral symmetry
breaking and no phase transition.
Instead, the moduli space is smoothly described
at all scales by gauge-invariant composite operators.
The infrared degrees of freedom are:
\begin{itemize}
\item $3 \times 3 = 9$ mesons: $M^i{}_j = Q^i\tilde{Q}_j$\,,
\item $1$ baryon: $B = \epsilon_{abc}\,Q^a Q^b Q^c$\,,
\item $1$ antibaryon: $\tilde{B} = \epsilon^{abc}\,
  \tilde{Q}_a\tilde{Q}_b\tilde{Q}_c$\,,
\end{itemize}
subject to the \emph{quantum-modified constraint}
\begin{equation}
  \det M - B\tilde{B} = \Lambda^6\,,
  \label{eq:constraint}
\end{equation}
where $\Lambda$ is the dynamical scale.
The constraint is enforced in the superpotential
by a Lagrange multiplier~$X$.

Each meson $M^i{}_j$ is a chiral superfield
containing a complex scalar and a Weyl fermion.
The fermionic components are called \emph{mesinos}:
they are the composite fermions of the confined theory,
analogous to constituent quarks in hadronic physics
but carrying no colour.

The mass deformation breaks $\UU(3)_L \times \UU(3)_R$
to its Cartan, ensuring $M$ is diagonal at the vacuum.
We deform the theory by UV mass terms motivated
by the field-energy argument:
\begin{equation}
  W = X\bigl(\det M - B\tilde{B} - \Lambda^6\bigr)
    + \sum_i m_i\,M^i{}_i\,,
  \label{eq:W}
\end{equation}
with $m_i = \mu\,q_i^2$, where $q_i = z_0 + z_i$ are
the family charges defined in Eq.~(\ref{eq:qi}).
On the mesonic branch ($B = \tilde{B} = 0$),
the $F$-flatness conditions
$\partial W/\partial M^i{}_i = 0$ and
$\partial W/\partial X = 0$ give the vacuum
(writing $P \equiv m_1 m_2 m_3$):
\begin{equation}
  \langle M_i \rangle = \frac{\Lambda^2\,P^{1/3}}{m_i}
  \;\propto\; \frac{1}{q_i^2}\,.
  \label{eq:Mvev}
\end{equation}
The key observation is that the VEVs are \emph{inversely}
proportional to the UV masses:
a heavy preon produces a light meson VEV, and vice versa.

The mesino mass matrix is obtained from
$\partial^2 W / \partial M_i \partial M_j$
evaluated at the vacuum.
The $F$-term condition
$\partial W/\partial M^i{}_i = X\,\adj(M)_i + m_i = 0$
gives $X = -m_i/\adj(M)_i$.
Since $\adj(M)_i = \Lambda^6/M_i$ and
$M_i = \Lambda^2 P^{1/3}/m_i$ (Eq.~\ref{eq:Mvev}),
one finds $\adj(M)_i = \Lambda^4 m_i/P^{1/3}$,
so $m_i/\adj(M)_i = P^{1/3}/\Lambda^4$
is generation-independent.
Therefore $X = -(P/\Lambda^{12})^{1/3}$.
Since $X$ is generation-independent,
the mass matrix
$\partial^2 W/\partial M_i\partial M_j$
is diagonal with entries
$m_i + X\,\partial\adj(M)_i/\partial M_i$.
Evaluating on the vacuum~\cite{Csaki:1996zb},
the diagonal mesinos acquire holomorphic masses
\begin{equation}
  m_{\tilde{M}_i} \;\propto\; m_i \;=\; \mu\,q_i^2\,,
  \label{eq:mesino_mass}
\end{equation}
where the proportionality constant is
generation-independent (it depends on $\Lambda$
and $P$ but not on~$i$) and therefore cancels
in $K$.
This is exact at the holomorphic level: by the
non-renormalisation theorem~\cite{Seiberg:1993vc},
the superpotential mass parameters $m_i$ receive no
perturbative corrections.
The physical (pole) masses of the mesinos are
$m_{\rm phys,\,i} = Z_i^{-1}\,m_i$,
where $Z_i$ is the K\"ahler metric for the
composite field~$M^i{}_i$.
For physical masses to satisfy $K = 2/3$,
one requires $Z_i$ to be generation-independent
at leading order.
This is an assumption whose justification
proceeds as follows.
Before the mass deformation, the microscopic theory
has an $\SU(3)_L \times \SU(3)_R$ global symmetry
under which all flavours are equivalent, so
$Z_i = Z$ identically.
The mass deformation $W \supset m_i M_i$ breaks this
symmetry; corrections to the K\"ahler metric are
constrained by dimensional analysis and reality to
enter as $\delta Z_i/Z \sim c\,(m_i/\Lambda)^2$,
where $c$ is an $\mathcal{O}(1)$ coefficient.
For the heaviest lepton ($m_\tau \approx 1.78\GeV$)
with $\Lambda \sim 1\TeV$,
$\delta Z_\tau/Z \sim 3 \times 10^{-6}$,
below the observed $0.001\% = 10^{-5}$ precision
of the lepton ratio.
For the inverse (quark) sector, the relevant
masses enter through the meson VEVs
$\langle M_i \rangle \propto 1/m_i$;
the same estimate gives
$\delta Z_b/Z \sim (m_b/\Lambda)^2 \sim 2 \times 10^{-5}$,
consistent with the coarser $0.3\%$ quark-sector
precision but only marginally so.
QCD corrections are an additional source of
generation-dependent K\"ahler dressing;
however, since QCD is flavour-blind
(the anomalous dimension $\gamma_m$ is
flavour-universal to five loops~\cite{Martin:2019}),
these enter only at order
$\alpha_s \times (m_i/\Lambda)^2$, which is subdominant.
A rigorous determination of~$Z_i$ requires the
non-perturbative K\"ahler potential,
which is beyond current techniques;
the estimates above are necessary-order checks,
not proofs of generation-universality.
The $q_i^2$ pattern is an input (set in the UV by the
field-energy argument); the virtue of the s-confining
model is that it \emph{preserves} this pattern to
the infrared and simultaneously generates the
inverse sector.

\section{Direct and inverse sectors}

The confined theory produces two distinct mass spectra
from a single meson matrix.
To see why, note that $M$ enters the physics in two ways.

\paragraph{Direct sector: mesino masses.}
The mesinos---the fermionic partners of the meson
scalars---have masses $m_i = \mu\,q_i^2$.
These are the UV mass parameters themselves, protected
by holomorphy.
Since $m_i \propto q_i^2$, the ratio
$K = \sum q_i^2/(\sum q_i)^2 = 2/3$ is automatic.

\paragraph{Inverse sector: meson VEVs.}
The scalar meson VEVs are
$\langle M_i \rangle \propto 1/m_i \propto 1/q_i^2$
(Eq.~\ref{eq:Mvev}).
If an additional field acquires its mass through a
flavour-diagonal Yukawa coupling to
$\langle M \rangle$, it will have
$m_{\rm phys} \propto \langle M_i \rangle \propto 1/q_i^2$.
The identification of such a field with the
down-type quarks is an additional input to the
framework, motivated by three observations:
(i)~the down-quark hierarchy
$m_d \ll m_s \ll m_b$ matches
the inverse-charge ordering
$1/q_1^2 \ll 1/q_2^2 \ll 1/q_3^2$;
(ii)~$K^{-1}(u,c,t) = 0.92$
deviates from $2/3$ by $38\%$,
ruling out the up quarks;
and (iii)~among the $\binom{6}{3} = 20$
quark triples, $(d,s,b)$ gives the smallest
$|K^{-1} - 2/3|$.
A first-principles derivation would require
specifying the SM embedding
(see ``Open sectors'' below);
the present letter treats the sector assignment
as phenomenological, with the statistical
significance addressed under ``Look-elsewhere.''
Denoting these inverse-sector masses by
$\hat m_i \propto 1/q_i^2$
(to distinguish them from the direct-sector
$m_i \propto q_i^2$), the \emph{inverse} ratio
(we write $K^{-1}$ as shorthand for
``$K$ evaluated on inverse masses,''
i.e.\ $K(1/\hat m_1,1/\hat m_2,1/\hat m_3)$,
not the reciprocal~$1/K$)
\begin{equation}
  K^{-1} \;\equiv\;
  K\!\left(\frac{1}{\hat m_1},\frac{1}{\hat m_2},
  \frac{1}{\hat m_3}\right)
  = \frac{\sum q_i^2}{\bigl(\sum q_i\bigr)^2}
  = \frac{2}{3}
  \label{eq:Kinv_derived}
\end{equation}
is equally automatic: since
$1/\hat m_i \propto q_i^2$,
the ratio reduces to the same normalisation identity.

This is the central structural result.
The same $\SU(3)$ gauge dynamics that confines the
preons into mesons automatically provides two sets
of composite operators---$M$ and $M^{-1}$---whose
mass spectra are related by inversion.
One set gives $K = 2/3$; the other gives $K^{-1} = 2/3$.
Both are consequences of the same underlying
charge normalisation, not separate numerical accidents.

\paragraph{The adjugate matrix.}
The relationship $M \leftrightarrow M^{-1}$ can be
stated invariantly using the adjugate (classical adjoint):
\begin{equation}
  \adj(M) \;=\; (\det M)\,M^{-1}
  \;=\; \Lambda^6\,M^{-1}\,,
  \label{eq:adjugate}
\end{equation}
where the last equality uses the quantum
constraint~(\ref{eq:constraint}).
The adjugate eigenvalues are
$\adj(M)_i = \Lambda^6/\langle M_i\rangle
\propto m_i \propto q_i^2$.
Thus $\adj(M)$ and $M$ carry complementary information:
\begin{equation}
  \langle M_i\rangle \propto 1/q_i^2\,,
  \qquad
  \adj(M)_i \propto q_i^2\,.
  \label{eq:dual_pair}
\end{equation}%
\footnote{The composite pair $(M_i, \adj(M)_i)$ has the
mathematical structure of an $\mathcal{N}{=}2$
hypermultiplet~\cite{SeibergWitten:1994a};
this doublet structure emerges from confinement
plus isospin and provides a natural pairing
of the direct and inverse branches.
A quantitative derivation of the inter-branch
coupling is left to future work.}

The $\SU(3)_F^3$ anomaly cancels
automatically ($\Nf$ fundamentals vs.\ $\Nf$ antifundamentals),
and 't~Hooft anomaly matching between UV quarks and
IR composites is satisfied for
$\Nf = \Nc$~\cite{Csaki:1996zb}.

%----------------------------------------------------------------------
\section{Experimental verification}
We now compare the
predictions of the framework with data.
The sector angle~$\alpha$ is a free parameter
(a modulus of the family-charge direction).

\paragraph{Direct sector (charged leptons).}
The two parameters $\alpha_{\rm lep}$ and $\mu$
are jointly determined from the two inputs
$m_e$ and $m_\mu$, giving
$\alpha_{\rm lep} = 0.2222\,\mathrm{rad}$ and
$\mu = 1883\MeV$.
The formula $m_i = \mu\,q_i^2$ then predicts
$m_\tau = 1776.97\MeV$, to be compared with
$m_\tau^{\rm exp} = 1776.86 \pm 0.12\MeV$~\cite{PDG:2024}
($0.9\sigma$; the model encodes a one-parameter
relation among three observables).
Equivalently, the experimental Koide ratio
\begin{equation}
  K(e,\mu,\tau) \;=\;
  \frac{\sum m_i}{\bigl(\sum \sqrt{m_i}\bigr)^2}
  \;=\; 0.666661\,,
  \label{eq:K_lep}
\end{equation}
deviates from the predicted value $2/3$
by $0.001\%$~\cite{PDG:2024}.
This ratio was first observed empirically by
Koide~\cite{Koide:1983qe} in a formulation using
square roots of masses; the charge-squared form
makes the gauge origin transparent.
The radiative stability of the relation was
studied by Sumino~\cite{Sumino:2009}
in the context of a gauged $\UU(3)$ family
symmetry~\cite{Koide:2006},
and independently by
G\'erard~\cite{Gerard:2022}.

\paragraph{Inverse sector (down quarks).}
Since $K$ is a ratio of mass sums, it is invariant
under a common rescaling $m_i \to c\,m_i$
(the factor~$c$ cancels between numerator
and denominator).
When the QCD anomalous dimension $\gamma_m$ is
flavour-universal---as it is to five-loop
order~\cite{Martin:2019}---all $\overline{\rm MS}$
masses at scale~$Q$ are related by such a common
rescaling, so $K^{-1}$ is scale-independent
to the accuracy of the flavour-universal
QCD anomalous dimension.
Using the PDG~2024 recommended
values~\cite{PDG:2024}
($m_d(2\GeV) = 4.70 \pm 0.07\MeV$,
$m_s(2\GeV) = 93.5 \pm 0.8\MeV$,
$m_b(m_b) = 4.183 \pm 0.007\GeV$),
the inverse ratio is
\begin{equation}
  K^{-1}(d,s,b) \;\equiv\;
  \frac{\sum_i m_i^{-1}}
  {\bigl(\sum_i m_i^{-1/2}\bigr)^2}
  \;=\; 0.6647\,,
  \label{eq:Kinv}
\end{equation}
a $0.3\%$ deviation from $2/3$.
(This is a conservative estimate from the
mixed-scale PDG convention; at a common scale
$Q = 2\GeV$, the deviation shrinks to $\sim 0.03\%$.)
The deviation is well within the $\pm 1\sigma$ error band
(Table~\ref{tab:K}) and is small compared with the
naive K\"ahler correction scale
$\alpha_s/\pi \sim 7\%$ at $Q = m_b$,
suggesting partial cancellation between
perturbative and non-perturbative
contributions~\cite{Koide:2005,Zenczykowski:2012,RodejohannZhang:2011}.

\begin{table}[b]
\caption{Predicted and observed ratios.
  Lepton masses: pole values~\cite{PDG:2024}.
  Quark masses: PDG~2024 recommended
  values~\cite{PDG:2024}
  ($m_d = 4.70 \pm 0.07$,
  $m_s = 93.5 \pm 0.8\MeV$ at $Q = 2\GeV$;
  $m_b(m_b) = 4.183 \pm 0.007\GeV$;
  the mixed-scale convention is justified
  by the flavour-universal QCD $\gamma_m$).
  Since $K$ is invariant under $m_i \to c\,m_i$
  and $\gamma_m$ is flavour-universal,
  $K^{-1}$ is scale-independent;
  Fig.~\ref{fig:running} confirms stability across
  14 decades.
  Uncertainties propagated from mass errors via
  numerical $\partial K / \partial m_i$.
  The lepton sector has two fitted parameters
  ($\alpha_{\rm lep}$, $\mu$); the genuine
  one-parameter prediction is $m_\tau^{\rm pred} =
  1776.97\MeV$ vs.\ $m_\tau^{\rm exp} =
  1776.86 \pm 0.12\MeV$~\cite{PDG:2024} ($0.9\sigma$).
  The quark sector is a zero-parameter phenomenological test.}
\label{tab:K}
{\small
\begin{tabular}{lcccc}
  \hline\hline
  Sector & Predicted & Observed & $1\sigma$ & Dev.\ ($n\sigma$) \\
  \hline
  $(e,\mu,\tau)$ direct    & $2/3$  & $0.666661$ & $0.000007$ & $0.9$ \\
  $(d,s,b)$ inverse        & $2/3$  & $0.6647$   & $0.0017$ & $1.2$   \\
  \hline\hline
\end{tabular}}
\end{table}

Rather than proposing an isolated mass relation,
the framework exhibits a nontrivial reduction
of fermion-sector structure:
the charged-lepton (direct) and inverse down-quark
tuples appear as two faces of a common
supersymmetric origin.

\paragraph{RG stability.}
Using the five-loop $\overline{\rm MS}$ running
implemented in SMDR~\cite{Martin:2019},
Fig.~\ref{fig:running} shows $K^{-1}(d,s,b)$
and $K(e,\mu,\tau)$ from $Q = 1\GeV$ to the Planck scale.
The $\pm 1\sigma$ bands (propagated from PDG mass uncertainties)
demonstrate that both ratios remain within $0.3\%$
of~$2/3$ across the entire range.
The inverse down-quark relation is consistent with
$K^{-1} = 2/3$ at every scale tested.

\begin{figure}[t]
\centering
\includegraphics[width=\columnwidth]{figures/koide_inv_bw.pdf}
\caption{Running of $K^{-1}(d,s,b)$ (solid) and $K(e,\mu,\tau)$
(dashed) from $Q = 1\GeV$ to the Planck scale,
computed with SMDR 5-loop RG evolution~\cite{Martin:2019}.
The purpose of the figure is to demonstrate
scale-independence, not to read off absolute values
(which should be taken from Table~\ref{tab:K},
since the SMDR default light-quark inputs differ from
PDG~2024 by up to~$7\%$).
Shaded: $\pm 1\sigma$ from PDG mass uncertainties.
Dotted: particle thresholds.}
\label{fig:running}
\end{figure}

%----------------------------------------------------------------------
\section{Discussion}
The overall scale
$\mu \approx 1883\MeV$ and the sector angle~$\alpha$
are not determined by the present argument.
Computing~$\mu$ requires the K\"ahler metric on the
quantum-modified moduli space, which is not constrained
by holomorphy.

\paragraph{Falsifiable bound.}
The radiative stability of $K = 2/3$ constrains the
family symmetry breaking scale.
Below~$\Lambda_F$, SM gauge loops (principally
electromagnetic) generate flavour-dependent
self-energy corrections that would spoil
$K = 2/3$ unless compensated by family gauge
boson exchange.
Sumino~\cite{Sumino:2009} showed that the leading
electromagnetic correction is cancelled by the
exchange of degenerate $\UU(3)_F$ gauge bosons,
with the residual correction suppressed by
family gauge boson mass splittings.
The observed $\delta K < 10^{-5}$ for charged leptons
then requires the family gauge boson masses
$M_F \lesssim \mathcal{O}(10)\TeV$.
In our framework, Higgs--confinement
complementarity~\cite{Csaki:1996zb}
ensures that the confining phase is smoothly
connected to a Higgsed phase in which the
$\SU(3)_F$ gauge bosons acquire masses
$M_F \sim g_F \langle M \rangle^{1/2}$
and can mediate the Sumino cancellation.
(Since the meson VEVs span a large dynamical
range, the resulting gauge boson masses are
non-degenerate; whether the residual corrections
preserve $\delta K < 10^{-5}$ requires a
quantitative analysis that we leave to future work.)
Non-observation of flavour-changing neutral currents
(e.g.\ $\mu \to e\gamma$, $K$--$\bar{K}$ mixing)
from family gauge boson exchange constrains
$M_F \gtrsim \mathcal{O}(1)\TeV$;
the window $1$--$10\TeV$ is testable,
contingent on the Sumino cancellation surviving
the non-degeneracy noted above.

\paragraph{Look-elsewhere.}
The framework predicts that if $K = 2/3$ holds
for a direct sector, there exists an inverse sector
satisfying $K^{-1} = 2/3$, but does not specify
which SM fermions populate each sector.
The direct branch (charged leptons) is selected
a~priori by the Koide relation
itself~\cite{Koide:1983qe}.
The inverse branch (down quarks) is a phenomenological
identification; its statistical weight must
account for the freedom to search among
alternative triples.
Since the inverse sector arises from meson VEVs
in the \emph{same} confined gauge theory
that produces the direct (lepton) sector,
only colour-carrying composites---i.e.\ quarks---are
candidates; cross-sector (lepton--quark) triples are
excluded by this structural constraint.
Among the $\binom{6}{3} = 20$ triples of the
six quarks, only $(d,s,b)$ satisfies
$|K^{-1} - 2/3| < 0.5\%$;
a Monte Carlo estimate using $10^7$ random mass
spectra (6~quark masses drawn log-uniformly on
$[2, 1.8 \times 10^5]\MeV$) gives a probability
$\sim 5 \times 10^{-3}$ of the best inverse
triple lying within~$0.3\%$.
If both sectors are taken as pre-selected,
the joint probability drops to~$< 10^{-6}$.

\paragraph{Open sectors.}
The up-quark sector and neutrino masses are not addressed;
both require additional structure
(a seesaw or Majorana mechanism for the latter).
The 9~mesinos $M^i{}_j$ transform as
$(\mathbf{3},\bar{\mathbf{3}})$ under the
$\SU(3)_L \times \SU(3)_R$ flavour symmetry.
If $\SU(2)_L$ is embedded in $\SU(3)_L$ as the
upper $2 \times 2$ block, the decomposition
$\mathbf{3} \to \mathbf{2} + \mathbf{1}$
provides a natural split into doublet and singlet
composites, from which the SM quantum numbers
can in principle be assigned.
The full Standard Model embedding, electroweak charges,
and the origin of the CKM matrix through non-holomorphic
K\"ahler dressing are deferred to a companion work.

%----------------------------------------------------------------------
\section{Summary}
We have shown that:
(i)~the classical field energy of any $\UU(1)$ gauge
theory scales as $U \propto q^2$, motivating
the mass ansatz $m \propto q^2$
which we hypothesise survives in
the confining regime;
(ii)~$n = 2$ is the unique monomial for which the ratio
$K = \sum q_i^2/(\sum q_i)^2$ is direction-independent;
(iii)~a minimal $\mathcal{N}{=}1$ $\SU(3)$ gauge theory
with three flavours realises both the direct and inverse
sectors through its meson and adjugate composites.
The quadratic charge dependence of the
holomorphic masses is exact
(protected by the non-renormalisation theorem);
extending $K = 2/3$ to physical masses
requires a generation-universal K\"ahler metric,
which is natural at leading order but receives
non-perturbative corrections.
The one-parameter lepton prediction gives
$m_\tau = 1776.97\MeV$ ($0.9\sigma$ from experiment);
the zero-parameter inverse down-quark
phenomenological test agrees at~$1.2\sigma$.
If the Sumino cancellation mechanism~\cite{Sumino:2009}
can be adapted to the non-degenerate gauge boson
spectrum of the complementary Higgs phase---which
requires further quantitative analysis---family
gauge boson masses would be bounded to
$1$--$10\TeV$.

\paragraph*{Acknowledgments.}
During the preparation of this work the author used Claude
(Anthropic) for algebraic verification, numerical checks,
and drafting assistance.
The author reviewed and edited all content and takes full
responsibility for the publication.

\begin{thebibliography}{99}

\bibitem{Georgi:1999}
  H.~Georgi,
  \textit{Lie Algebras in Particle Physics},
  2nd ed.\ (Westview, 1999).

\bibitem{Foot:1994}
  R.~Foot,
  arXiv:hep-ph/9402242 (1994).

\bibitem{Seiberg:1994bz}
  N.~Seiberg,
  Nucl.\ Phys.\ B \textbf{435}, 129 (1995).

\bibitem{Csaki:1996zb}
  C.~Cs\'aki, M.~Schmaltz, and W.~Skiba,
  Phys.\ Rev.\ Lett.\ \textbf{78}, 799 (1997).

\bibitem{Seiberg:1993vc}
  N.~Seiberg,
  Phys.\ Lett.\ B \textbf{318}, 469 (1993).

\bibitem{SeibergWitten:1994a}
  N.~Seiberg and E.~Witten,
  Nucl.\ Phys.\ B \textbf{426}, 19 (1994);
  \textbf{431}, 484 (1994).

\bibitem{PDG:2024}
  S.~Navas \textit{et al.} (Particle Data Group),
  Phys.\ Rev.\ D \textbf{110}, 030001 (2024).

\bibitem{Koide:1983qe}
  Y.~Koide,
  Phys.\ Lett.\ B \textbf{120}, 161 (1983).

\bibitem{Zenczykowski:2012}
  P.~Zenczykowski,
  Phys.\ Rev.\ D \textbf{86}, 117303 (2012).

\bibitem{Sumino:2009}
  Y.~Sumino,
  JHEP \textbf{0905}, 075 (2009).

\bibitem{Gerard:2022}
  J.-M.~G\'erard,
  arXiv:2208.00364 (2022).

\bibitem{Koide:2005}
  Y.~Koide,
  Phys.\ Lett.\ B \textbf{595}, 163 (2004).

\bibitem{Koide:2006}
  Y.~Koide,
  Phys.\ Rev.\ D \textbf{73}, 057901 (2006).

\bibitem{Martin:2019}
  S.P.~Martin and D.G.~Robertson,
  Phys.\ Rev.\ D \textbf{100}, 073004 (2019).

\bibitem{FroggattNielsen:1979}
  C.D.~Froggatt and H.B.~Nielsen,
  Nucl.\ Phys.\ B \textbf{147}, 277 (1979).

\bibitem{IntriligatorSeibergShenker:1995}
  K.~Intriligator, N.~Seiberg, and S.H.~Shenker,
  Phys.\ Lett.\ B \textbf{342}, 152 (1995).

\bibitem{RodejohannZhang:2011}
  W.~Rodejohann and H.~Zhang,
  Phys.\ Rev.\ D \textbf{83}, 073012 (2011).

\bibitem{Rivero:2005}
  A.~Rivero,
  arXiv:hep-ph/0505220 (2005).

\bibitem{Ma:2006}
  E.~Ma,
  Phys.\ Lett.\ B \textbf{649}, 287 (2007).

\end{thebibliography}

\end{document}
